\section{Speech Synthesis}
\label{sec:tts}

As a final downstream task we consider speech synthesis or text-to-speech (TTS) where models output speech for a corresponding input text.
Most prior work focuses on English using clean corpora and high quality phonemizers~\citep{tan2021survey}.
These resources are only available for a small number of languages which is why expanding TTS to \nummmslang{} languages requires different choices.
\citet{black2019cmu} built TTS systems for 699 languages using data from a similar source as \MMS{}.
Encouraged by our initial results indicating higher quality data (\textsection\ref{sec:data-comp-cmu}), we extend their work by building models for \nummmslang{} languages.

We first describe the model architecture on which we build (\textsection\ref{sec:tts-model}), 
how we pre-process the data to train TTS models (\textsection\ref{sec:tts-preprocess}) and how we evaluate our models (\textsection\ref{sec:tts-eval}).
Next, we present an ablation of our design choices compared to a highly optimized setup used for English (\textsection\ref{sec:tts-results-english}).
We also measure performance when synthesizing out-of-domain data for 61 languages of the FLEURS benchmark (\textsection\ref{sec:tts-results-fleurs}), and finally present results for all \nummmslang{} languages (\textsection\ref{sec:tts-results-all}).



\subsection{Text-To-Speech Model}
\label{sec:tts-model}

Our Text-to-Speech (TTS) model is based on VITS~\citep{kim2021conditional}, which is one of the state-of-the-art TTS approaches. 
While VITS has previously been applied in a multilingual setting for English, Portuguese, and French~\citep{casanova2022yourtts}, we scale it to 1,107 languages.

VITS is an end-to-end speech generation network that predicts the raw speech waveform from a text sequence. 
It can be viewed as a conditional variational auto-encoder (VAE;~\citealt{Kingma2013AutoEncodingVB}) that estimates audio features based on the input text. The spectrogram-based acoustic features are generated by a flow-based sub-network, which includes multiple affine coupling layers~\citep{dinh2017density} and a text encoder. 
The waveform is decoded using a stack of transposed convolutional layers that have been adapted from HiFi-GAN~\citep{kong2020hifi}. The model is trained end-to-end with a combination of losses derived from variational lower bound and adversarial training. During inference, the text encodings are upsampled based on an internal duration prediction module and then mapped into the waveform using a cascade of the flow module and HiFi-GAN decoder.

We train separate VITS models for each language. 
Most of our hyperparameters are identical to the VITS model trained on LJSpeech~\citep{ljspeech17,kim2021conditional} except for these differences:
Instead of training the model for 800K steps, we train each model for 100K steps using eight V100-GPUs with a batch size of 64 per GPU.
This setup was only slightly worse than the original configuration but reduced training time by approximately eight times which made training a large number of TTS systems feasible (\textsection\ref{sec:tts-eval}).
We experimented with different learning rate settings and found that the original learning rate schedule of VITS worked best.


\subsection{Text and Speech Data Pre-processing}
\label{sec:tts-preprocess}

\paragraph{Data Selection.}
To train models we use the \MMS{} dataset which provides paired speech and text data.
For most languages we have a single recording of the New Testament, however, for 99 languages multiple recordings are available (\textsection\ref{sec:data-paired-source}). 
For these languages, we choose a single recording in order to avoid introducing additional speakers into the training data.
To choose a recording, we train ASR models on the data and select the recording based on which the corresponding ASR model achieves the lowest CER on a held-out set of an out-of-domain evaluation set. 
If no out-of-domain evaluation set is available, we choose a random recording. 
Finally, if there are both drama and non-drama recordings, then we consider only non-drama recordings.


\paragraph{Text Representations.}
Most current TTS models convert text to phonemes using grapheme-to-phoneme tools such as g2p~\citep{g2pE2019} which rely on a lexicon to map input characters to phonetic representations.
However, such lexicons are not available for most low resource languages since they require manual annotation.
In order to scale TTS to over one thousand languages, we represent the input text as individual letters for languages with a small vocabulary.
For languages with 200 or more characters, we use a uroman encoding (\citealt{hermjakob-etal-2018-box}; \textsection\ref{sec:data-initial-align}).
We validated this strategy on the languages of FLEURS where found that letter-based models outperform uroman-based systems across all languages, except for Amharic and Korean, which both have large character sets of between 200-1,000 characters.
We therefore used this mixed strategy.\footnote{
For all \nummmslang{} languages, the following languages use a uroman encoding: Amharic, Gumuz, Korean, Sebat Bet Gurage and Tigrinya.} 


\paragraph{Speech Data Pre-processing.}
For drama recordings, we remove background music to enhance the quality of TTS models.
We use a denoiser model~\citep{defossez2020real} to remove background music.
We also noticed that some utterances contain multiple speakers, usually voicing different characters in the read stories. 
We use a simple heuristic to detect those utterances and remove them from the training data.
The heuristic computes the variance of the pitch on voiced frames and removes utterances with high variance in the pitch; the pitch estimation is based on~\citet{Mauch2014PYINAF}.
We found that removing 15\% of the utterances with the highest pitch variance in each recording removed many multi-speaker utterances.



\subsection{Evaluation Methodology}
\label{sec:tts-eval}

Evaluation of speech synthesis is not straightforward even in the most researched English setting comprising a clean corpus with single speaker data~\citep{ljspeech17} and expanding evaluation to a large number of languages poses additional challenges.
Similar to prior work, we rely on both automatic metrics (MCD and ASR) and human studies which we detail next.

\paragraph{Mel-Cepstral Distortion (MCD).}
Mel-cepstral distortion is an automatic metric which measures the closeness of synthesized speech to a human utterance for the same text in terms of the warping distance of mel frequency cepstral coefficients.
There are two disadvantages of MCD:
first, it is only meaningful when the voice and prosody of the speaker as well as recording conditions in the training data matches the evaluation data because the mel cepstral sequences are sensitive to these aspects.
This prevents evaluation on data outside the \MMS{} domain.
Second, it is not well suited to measuring the intelligibility of the TTS output.
We address the former by focusing MCD evaluation on \MMS{} data and the latter via the next metric.

\paragraph{Automatic speech recognition (ASR).}
Transcribing the synthesized speech with automatic speech recognition and then measuring the error rate has the potential to address both issues of MCD: it can be measured on evaluation sets which are out-of-domain and it captures the content of the synthesized speech.
Specifically, we synthesize both in-domain data from the \MMS{} evaluation sets as well as out-of-domain data from the FLEURS corpus. 
Next, we apply ASR models and compute CER of the ASR model output with respect to the input text of the TTS system.
Low CER indicates that the TTS model is able to capture the content of the input text.
Unless otherwise stated, we use ASR models trained on FLEURS data.

\paragraph{Mean Opinion Score (MOS).}
Finally, we evaluate models using MOS and because it is very hard to find human raters proficient for a large number of languages, we ask raters to judge the fidelity of the synthesized samples together with how natural the speech sounds.
We rely on ASR to get a sense of how much of the content is preserved in the TTS outputs and consider the MOS score for generation quality and naturalness.
We evaluate 50 samples per method of each language, collect ten judgements per sample and ask raters to judge the quality from 1 (lowest quality) to 5 (highest quality). 
Results are reported in terms of the mean score as well as the 95\% confidence interval.
We use the CrowdMOS~\citep{ribeiro2011crowdmos} package with the recommended recipes for detecting and discarding inaccurate ratings.
We do not require raters to be able to speak the respective language (except for English) as it is very hard to find raters that speak all the languages we consider. 
We rely on the ASR error rate to get a sense of content preservation in the TTS models.


\subsection{Evaluation of Design Choices}
\label{sec:tts-results-english}

\subsubsection{Training Setup}

We first analyze the quality of our training setup (\textsection\ref{sec:tts-model}) and compare it to the common VITS setup~\citep{kim2021conditional} which trains models for up to 800K updates on the clean LJSpeech corpus~\citep{ljspeech17} with text data pre-processed by a high-quality phonemizer~\citep{g2pE2019}. 
This contrasts to our setup where we train the same model for 100K updates on the \MMS{} data using letters.

\insertTableResultsTTSDesign
\insertTableResultsTTSDrama


We evaluate performance on the development sets of \MMS{} which is in-domain for \MMS{} models, LJSpeech (LJS; \citealt{ljspeech17}) which is in-domain for the original VITS setup and FLEURS which is out-of-domain for both.\footnote{For measuring CER on LJSpeech, we remove punctuation and capitalization of both references and hypothesis using the normalization of~\citet{radford2022whisper}.}
The audio of FLEURS frequently contains high levels of noise and reverberation and it is relatively uncommon to use it for TTS, however, it does cover a large number of languages and it is out-of-domain with respect to the \MMS{} training data which makes it an interesting evaluation set, particularly for the experiment in \textsection\ref{sec:tts-results-fleurs}.

\autoref{tab:tts-design} shows that our reduced setup (row 5) generally performs less well than the highly optimized setup of VITS for LJSpeech~(row 2; \citealt{kim2021conditional}) and each design choice (fewer training updates, \MMS{} training data and character inputs) leads to a reduction in quality.
In terms of character error rate, the degradation is most pronounced on out-of-domain settings with respect to our reduced setup (row 5) and less so on the in-dmain \MMS{} development set.
We stress that these choices enable scaling TTS to over 1,000 languages at manageable compute requirements and no need for language-specific text processing tools.



Note that the CER of almost all models on FLEURS data is lower than for the corresponding natural speech and that the MOS scores of the human reference utterances is also lower than for other datasets.
We attribute this to the high levels of noise and reverberation in the FLEURS audio which results in increased CER for the original samples compared to the CER of the synthesized samples.

\subsubsection{Data with Background Music}

Next, we ablate how we build speech synthesis models based on recordings with background music which is a setting that applies to about 38\% of the languages (\textsection\ref{sec:data-paired-source}).
For this purpose, we train a model on an English recording with background music before and after the pre-processing steps outlined in~\textsection\ref{sec:tts-preprocess} and compare this to a model trained on another English recording that does not contain any background music to start with (no background music).
The pre-processing denoises the data and removes utterances with multiple speakers.

The results (\autoref{tab:tts-drama}) show that both denoising and removing samples with multi-speaker utterances results in performance improvements compared to the original data with background music.
Both steps reduce the CER gap to models trained on no background music data by 69-87\% relative to the CER of the system trained on data with background music.
MOS scores also increase after pre-processing, sometimes close to the level of the data with no background music on \MMS{} evaluation data (3.47 vs. 3.51).



\subsection{Out-of-Domain Evaluation}
\label{sec:tts-results-fleurs}

The \MMS{} data is from a particular narrow domain (\textsection\ref{sec:data-paired-source}) which poses the question of whether TTS models trained on this data will generalize to other domains.
To get a better sense of this, we train speech synthesis models and evaluate their quality on both in-domain and out-of-domain data.
As in-domain data we use the test sets of \MMS{} and as out-of-domain data we use FLEURS. 
This enables evaluation on 61 languages of the FLEURS benchmark (FLEURS-61) which are covered by \MMS{} data. 
However, a downside of FLEURS is that it contains high levels of noise which makes it challenging when comparing synthesized audio to the human reference audio.

\insertTableResultsTTSFLEURS

\autoref{tab:tts-fleurs} shows that \modelname{} models are robust to domain shift:
the CER of synthesized speech (TTS) is only slightly higher out-of-domain compared to in-domain and MOS scores for the synthesized samples are nearly identical.
We note that the in-domain and out-of-domain settings are measured on different data which does not enable strong claims about identical performance.
The systems also retain much of the original content as the small difference in CER between TTS and human utterances (ref) shows.

The MOS scores also indicate that our systems have lower sound quality compared to human utterances but the difference is not very large on in-domain data (3.51 vs. 3.61).
Unfortunately, out-of-domain MOS scores for the references are affected by the noisy speech in the FLEURS audio as noted earlier.
We conclude that TTS models trained on \MMS{} data perform well out-of-domain.


\subsection{Evaluation on \nummmslang{} Languages}
\label{sec:tts-results-all}

Finally, we train models for all languages of \MMS{} and focus on in-domain evaluation since finding out-of-domain evaluation data for such a large number of languages is difficult.
Specifically, we measure MCD and ASR CER on the \MMS{} test sets. 
To be able to evaluate ASR quality on all languages, we use ASR models trained on \MMS{} data which results in much lower error rates.
Similar to \textsection\ref{sec:asr-full-eval}, we group results into six geographical regions covered by \MMS{}: Asia, North America, South America, Europe, Africa and the Pacific region.

To get an overall sense of how many models are of good quality, we measure whether the model for a particular language has ASR CER $\leq$ 5. 
This indicates the number of systems which make on average no more than one error in twenty characters.
While this measure is by far not perfect, it enables us to get a broad sense of quality across a large number of languages.

\insertTableResultsTTSFull

\autoref{tab:tts-full} shows that about 85\% of the \nummmslang{} languages meet the CER quality threshold. 
South American and European languages achieve the highest rate at 95\% and African languages the lowest rate of 76\%. 
This is in part driven by different writing scripts.
The ASR character error rates are generally low because the error rates are based on ASR models trained on \MMS{} data.
